% from ega3-6-fr.tex
\oldpage[III]{153}

\[
  ''\mathscr E_{pq}^2
  =\bigoplus_{q_1+q_2+\cdots+q_m=q}
   \mathscr{T}\!or_p^S
   \bigl(\mathscr H_{q_1}(\mathscr P_\bullet^{(1)}),\ldots,
         \mathscr H_{q_m}(\mathscr P_\bullet^{(m)})\bigr).
  \tag{6.5.15.1}\label{III.6.5.15.1.fr}
\]
Nous étudierons dans (6.8) les suites spectrales d'associativité auxquelles
donnent lieu ces foncteurs hypertor d'un nombre quelconque de complexes.
\end{env}

\subsection{Foncteurs hypertor globaux de complexes de Modules quasi-cohérents
et suites spectrales de Künneth : cas de la base affine}
\label{subsection:III.6.6.fr}

\begin{env}[6.6.1]
Considérons un schéma affine $S=\operatorname{Spec}(A)$ et deux $S$-schémas
quasi-compacts $X^{(i)}$ ($i=1,2$); soit $\mathscr P_\bullet^{(i)}$ un complexe
de $\mathscr O_{X^{(i)}}$-Modules quasi-cohérents, limité inférieurement, dont
l'opérateur de dérivation est de degré $-1$ ($i=1,2$). Considérons d'autre
part un recouvrement fini
$\mathfrak U^{(i)}=(U_\alpha^{(i)})$ de $X^{(i)}$ par des ouverts affines;
soit
\[
  X=X^{(1)}\times_S X^{(2)},
\]
qui est un $S$-schéma quasi-compact (I, 5.5.1 et 6.6.4), et soit
$\mathfrak U=\mathfrak U^{(1)}\times_S\mathfrak U^{(2)}$ le recouvrement de
$X$ formé des ouverts affines
$U_\alpha^{(1)}\times_S U_\beta^{(2)}$. Pour tout couple d'entiers
$p\leqslant0$, $q\in\mathbf Z$, le groupe des $(-p)$-cochaînes alternées
\[
  C^{-p}(\mathfrak U^{(i)},\mathscr P_q^{(i)})
\]
du recouvrement $\mathfrak U^{(i)}$ à coefficients dans le faisceau
$\mathscr P_q^{(i)}$ (G, II, 5.1) est un $A$-module; pour $p>0$, on posera
$C^p(\mathfrak U^{(i)},\mathscr P_q^{(i)})=0$; on a ainsi défini un
bicomplexe
\[
  C^\bullet(\mathfrak U^{(i)},\mathscr P_\bullet^{(i)})
\]
de $A$-modules, dont les deux opérateurs de dérivation sont de degré $-1$. Il
résulte des définitions ($0_{\mathrm I}$, 12.1.2) que le $A$-module
d'homologie
\[
  \mathrm H_n\bigl(C^\bullet
  (\mathfrak U^{(i)},\mathscr P_\bullet^{(i)})\bigr)
\]
de ce bicomplexe (considéré à l'ordinaire comme un complexe simple pour le
degré total) n'est autre que le $A$-module d'hypercohomologie
\[
  \mathbf H^{-n}
  (\mathfrak U^{(i)},\mathscr P^{(i)\bullet}),
\]
où $\mathscr P^{(i)\bullet}$ est le complexe à opérateur de dérivation de
degré $+1$ obtenu en prenant $\mathscr P_{-q}^{(i)}$ pour composante de degré
$q$; par abus de notation, nous l'écrirons
$\mathbf H^{-n}(\mathfrak U^{(i)},\mathscr P_\bullet^{(i)})$. Il résulte alors
de (6.2.2) que ce $A$-module est canoniquement isomorphe au $A$-module
d'hypercohomologie
$\mathbf H^{-n}(X^{(i)},\mathscr P^{(i)\bullet})$, que nous écrirons de même
$\mathbf H^{-n}(X^{(i)},\mathscr P_\bullet^{(i)})$; il ne dépend donc pas du
recouvrement fini $\mathfrak U^{(i)}$ choisi.
\end{env}

\begin{env}[6.6.2]
Nous allons appliquer aux deux bicomplexes de $A$-modules
\[
  L_{\bullet\bullet}^{(i)}
  =C^\bullet(\mathfrak U^{(i)},\mathscr P_\bullet^{(i)})
  \qquad(i=1,2)
\]
et au bifoncteur covariant
$L_{\bullet\bullet}^{(1)}\otimes_A L_{\bullet\bullet}^{(2)}$ en ces deux
bicomplexes, la théorie générale de l'hyperhomologie des foncteurs par rapport
aux bicomplexes ($0_{\mathrm I}$, 11.7.4). Comme les cochaînes considérées sont
alternées et les recouvrements $\mathfrak U^{(i)}$ finis, on notera que les
modules $C^{-p}(\mathfrak U^{(i)},\mathscr P_q^{(i)})$ ne sont $\ne0$ que pour
un nombre fini (indépendant de $q$) de valeurs de $p$, et en particulier les
deux degrés de chacun des $L_{\bullet\bullet}^{(i)}$ sont limités
inférieurement. Nous désignerons par
\[
  \mathbf{Tor}_n^A(L_{\bullet\bullet}^{(1)},L_{\bullet\bullet}^{(2)})
\]
ou
\[
  \mathbf{Tor}_n^S
  (\mathfrak U^{(1)},\mathfrak U^{(2)};
   \mathscr P_\bullet^{(1)},\mathscr P_\bullet^{(2)})
\]
le $n$-ième module d'hyperhomologie de
$L_{\bullet\bullet}^{(1)}\otimes_A L_{\bullet\bullet}^{(2)}$, que nous
appellerons l'\emph{hypertor} d'indice $n$ de
$\mathscr P_\bullet^{(1)}$ et $\mathscr P_\bullet^{(2)}$, relatif aux
recouvrements $\mathfrak U^{(1)}$ et $\mathfrak U^{(2)}$. Lorsque
$\mathscr P_\bullet^{(1)}$ et $\mathscr P_\bullet^{(2)}$ sont réduits à leurs
termes de degré 0, $\mathscr F^{(1)}$ et $\mathscr F^{(2)}$, on écrit
\[
  \mathbf{Tor}_n^S
  (\mathfrak U^{(1)},\mathfrak U^{(2)};
   \mathscr F^{(1)},\mathscr F^{(2)})
\]
leur hypertor. On désignera par
\[
  \mathbf{Tor}_n^S
  (\mathfrak U^{(1)},\mathfrak U^{(2)};
   \mathscr P_\bullet^{(1)},\mathscr P_\bullet^{(2)}),
\]
suivant les conventions générales, le bicomplexe dont le composant d'indices
$(j,k)$ est
\[
  \mathbf{Tor}_n^S
  (\mathfrak U^{(1)},\mathfrak U^{(2)};
   \mathscr P_j^{(1)},\mathscr P_k^{(2)}).
\]

Comme $L_{\bullet\bullet}^{(i)}$ est un foncteur exact en
$\mathscr P_\bullet^{(i)}$, puisque les intersections des ensembles
